\begin{exercise}{130}
Show that for an $\infty$-groupoid $X$, there is a canonical Joyal equivalence
$X^{\mathrm{op}}\simeq X$.

\begin{proof}
Reverse barycentric coordinates by
$R_n(t_0,\ldots,t_n)=(t_n,\ldots,t_0)$.
For $\alpha\colon[n]\to[m]$, put
$\bar\alpha(i)=m-\alpha(n-i)$, as in Exercise~\ref{ex:21}.
The affine maps on topological simplices satisfy
\[
 R_m\alpha_*=(\bar\alpha)_*R_n:
 \quad\text{both $j$th coordinates are }
 \sum_{\alpha(i)=m-j}t_i.
\]
Thus $[\sigma,t]\mapsto[\sigma,R_nt]$ respects the realization relations
and defines a natural homeomorphism
$r_X\colon|X^{\mathrm{op}}|\to|X|$, whose inverse is the same construction
for $X^{\mathrm{op}}$.

By Exercise~\ref{ex:21}, $X^{\mathrm{op}}$ is Kan.
The fact granted in Exercise~\ref{ex:97} makes the unit maps
$\eta_X$ and $\eta_{X^{\mathrm{op}}}$ homotopy equivalences.
Hence the natural zigzag
\[
 X^{\mathrm{op}}\xrightarrow{\eta_{X^{\mathrm{op}}}}
 \mathcal S|X^{\mathrm{op}}|
 \xrightarrow{\mathcal S(r_X)}\mathcal S|X|
 \xleftarrow{\eta_X}X
\]
consists of homotopy equivalences between Kan complexes
(Lemma~1.1.47), hence of Joyal equivalences by Observation~2.2.11.
Choosing a homotopy inverse $q_X$ of $\eta_X$ gives the representative
$q_X\mathcal S(r_X)\eta_{X^{\mathrm{op}}}\colon X^{\mathrm{op}}\to X$.
Its homotopy class is independent of that choice, since the inverse of
$[\eta_X]$ in the homotopy category is unique.
Naturality of the zigzag makes this class canonical and natural in $X$.
\end{proof}
\end{exercise}
